\documentclass[12pt]{article}
\usepackage[utf8]{inputenc}
\usepackage[T1]{fontenc}
\usepackage{hyperref}
\title{Pannei Language Resources for K-12 Curriculum Development}
\author{Compiled by Claude Sonnet 5.5 (claude-sonnet-5-5) for Andr\'as Kornai}
\date{}
\begin{document}
\maketitle

\section{Introduction}
Pannei (ISO 639-3 \texttt{pnc}, Glottocode pann1239) is an Austronesian language of Sulawesi, Indonesia. Sources differ on its location: Wikipedia gives South Sulawesi, while Joshua Project and Asia Harvest place the people in West Sulawesi (Polewali Mamasa area). Two dialects, Tapango and Bulo, are reported. Speaker counts range from about 9,000 (1983) to about 12,000 (Ethnologue 2016, via Joshua Project). Almost all Pannei people are Muslim. The search found no school materials, primers, dictionaries or audio in Pannei; only encyclopedic and people-group profile pages.

\section{Primary Educational Resources}
No primers, textbooks or school programmes in Pannei were found in the search results.

\section{Linguistic and Reference Materials}
\begin{itemize}
\item Wikipedia, ``Pannei language'': overview of classification, dialects and speaker numbers. \url{https://en.wikipedia.org/wiki/Pannei_language}
\end{itemize}

\section{Religious and Community Resources}
\begin{itemize}
\item Joshua Project, Pannei in Indonesia: people-group profile; lists Scripture status as ``Translation Started''. \url{https://joshuaproject.net/people_groups/14301}
\item Asia Harvest, Pannei profile (PDF). \url{https://people-groups.asiaharvest.org/Indonesia/Pannei.pdf}
\item Praying for Indonesia, page on the Pannei. \url{https://www.prayingforindonesia.com/?p=843}
\end{itemize}

\section{Digital and Audio Resources}
No audio, video, corpus or software resources for Pannei were found.

\section{Curriculum Development Context}
Pannei is not known to be taught in schools, and no written educational material was located. Any curriculum would need to start from primary documentation (word lists, orthography, recorded texts) with community partners. Related Pitu Ulunna Salu languages may offer comparative material, but none was found in this search.

\end{document}